\chapter{Prova Discursiva — Peça Técnica e Situações-Problema}

\section{Formato e estratégia}

Para Auditor Estadual de Controle Externo, a prova discursiva vale 40 pontos e é composta por uma \textbf{peça de natureza técnica de até 60 linhas}, valendo 20 pontos, e \textbf{duas questões de situação-problema de até 30 linhas}, valendo 10 pontos cada. O treino deve reproduzir o limite físico e o tempo de prova.

\begin{teoria}[princípio central]
A discursiva do Cebraspe não premia texto ornamental. Ela premia \textbf{atendimento ao comando + conteúdo técnico correto + organização + precisão vocabular}. Antes de escrever, converta o comando em uma lista de itens obrigatórios e distribua linhas para cada um.
\end{teoria}

\section{Método C-A-E-R para peça técnica}

Use quatro blocos mentais:
\begin{enumerate}
\item \textbf{C — Contexto/critério}: qual norma, framework, requisito ou objetivo deveria ser atendido?
\item \textbf{A — Achado/análise}: o que foi observado e por que é inadequado?
\item \textbf{E — Efeito/risco/evidência}: que consequência pode ocorrer e que evidências sustentam a conclusão?
\item \textbf{R — Recomendação}: qual ação concreta, proporcional e verificável deve ser proposta?
\end{enumerate}

Esse formato não precisa aparecer como título na prova, mas serve como arquitetura de raciocínio.

\section{Como elaborar um achado de auditoria}

Um achado robusto conecta:
\[
\text{critério} \rightarrow \text{condição} \rightarrow \text{causa} \rightarrow \text{efeito/risco} \rightarrow \text{evidência} \rightarrow \text{encaminhamento}
\]

\textbf{Critério} é o padrão esperado. \textbf{Condição} é o fato constatado. \textbf{Causa} explica por que ocorreu. \textbf{Efeito} é consequência observada ou risco plausível. \textbf{Evidência} é suporte suficiente e apropriado. \textbf{Encaminhamento} deve atacar causa e risco, e não apenas repetir o problema.

\section{Distribuição das 60 linhas}

Uma estratégia segura para peça de 60 linhas:
\begin{itemize}
\item 4--6 linhas: enquadramento e síntese do problema;
\item 35--42 linhas: desenvolvimento dos pontos exigidos, com análise técnica;
\item 8--12 linhas: recomendações e conclusão;
\item margem mental restante: correções e transições.
\end{itemize}

Não existe obrigação de usar essa proporção; ela evita gastar metade da folha na introdução.

\section{Questões de 30 linhas}

Para questões situacionais, adote estrutura de três blocos:
\begin{enumerate}
\item resposta direta à tese do comando;
\item desenvolvimento dos subitens, um por parágrafo ou período organizado;
\item conclusão técnica com consequência/recomendação quando solicitada.
\end{enumerate}

\section{Erros que derrubam nota}

\begin{atencao}
Evite: responder tema diferente do comando; não tratar um dos subitens; listar siglas sem explicar aplicação; recomendar tecnologia sem conectar ao risco; confundir evidência com opinião; usar linguagem absoluta onde há trade-off; ultrapassar limite de linhas; e escrever conclusões genéricas como ``deve-se melhorar a segurança''.
\end{atencao}

\section{Treino 1 — Peça técnica: contrato SaaS e lock-in}

\begin{discursiva}[peça técnica — até 60 linhas]
Em auditoria de contrato de plataforma SaaS crítica, foram observados: ausência de cláusula de portabilidade; exportação disponível apenas em formato proprietário; contas privilegiadas compartilhadas entre técnicos da contratada; logs administrativos mantidos exclusivamente pelo fornecedor; indisponibilidade de 99,4\% em mês cujo SLA exige 99,9\%; e pagamentos efetuados com base em relatório produzido pela própria contratada, sem validação independente. Elabore peça técnica contendo: (a) principais achados e riscos; (b) evidências adicionais a obter; (c) relação entre medição e pagamento; e (d) recomendações prioritárias.
\end{discursiva}

\begin{resumo}[espelho de correção]
\textbf{Elementos esperados:} lock-in e ausência de estratégia de saída; risco de confidencialidade e responsabilização pelo compartilhamento de contas; baixa confiabilidade de logs sob domínio exclusivo da contratada; descumprimento do SLA; fragilidade de evidência para liquidação/pagamento; necessidade de obter contrato/TR/ETP, trilhas de acesso, logs independentes, cálculo do SLA, histórico de incidentes e pagamentos; recomendar portabilidade, formatos abertos, plano de saída testado, PAM/identidades individuais, envio imutável de logs para domínio do órgão, medição independente e glosa/ajuste conforme regras contratuais.
\end{resumo}

\begin{teoria}[resposta-modelo condensada]
A situação revela deficiências de planejamento, segurança, fiscalização e mensuração contratual. A ausência de requisitos de portabilidade e o uso exclusivo de formato proprietário elevam o risco de dependência do fornecedor e dificultam transição ao término do ajuste. O compartilhamento de contas privilegiadas elimina responsabilização individual e amplia risco de uso indevido; adicionalmente, a guarda exclusiva dos logs pela contratada reduz independência e confiabilidade da evidência.

A disponibilidade de 99,4\% deve ser confrontada com a fórmula, janela e exclusões definidas no SLA de 99,9\%. Se confirmado o descumprimento, a medição deve produzir os efeitos previstos no contrato, inclusive glosa ou sanção quando cabível. Pagamento sustentado apenas por relatório da parte interessada é controle frágil.

Devem ser obtidos ETP, termo de referência, contrato e aditivos, metodologia de medição, logs brutos, registros do monitoramento do órgão, chamados e incidentes, trilhas de contas privilegiadas e memórias de cálculo dos pagamentos. Recomenda-se instituir estratégia de saída e teste periódico de exportação, formatos interoperáveis, contas privilegiadas nominativas e gerenciadas, centralização segura de logs sob domínio do órgão, monitoramento independente e vinculação objetiva do pagamento ao resultado efetivamente medido.
\end{teoria}

\section{Treino 2 — Peça técnica: ransomware e continuidade}

\begin{discursiva}[peça técnica — até 60 linhas]
Um órgão foi atingido por ransomware. Os servidores estavam em RAID 6, havia snapshots diários no mesmo storage, o backup externo mais recente tinha 18 dias, contas administrativas não utilizavam MFA e os testes de restauração não eram realizados havia dois anos. A direção afirma que ``RAID e snapshots eram suficientes''. Produza peça técnica abordando causas de fragilidade, impacto sobre RPO/RTO, evidências a coletar e recomendações.
\end{discursiva}

\begin{resumo}[espelho de correção]
Distinguir disponibilidade de disco, snapshot e backup; explicar domínio de falha comum e ransomware; calcular que backup de 18 dias pode implicar perda potencial muito superior a RPO típico; ausência de teste torna RTO não demonstrado; evidências: política, inventário, logs, cópias, configuração, imutabilidade, restore, IAM, timeline; recomendações: 3-2-1, cópia offline/imutável, segmentação, MFA/PAM, EDR, testes, BIA, RPO/RTO formalizados, exercícios de resposta e recuperação.
\end{resumo}

\section{Treino 3 — Peça técnica: engenharia de dados}

\begin{discursiva}[peça técnica — até 60 linhas]
O TCE recebe dados mensais de folha de pagamento dos jurisdicionados. A equipe descobre duplicidades, CPF inválido, valores retroativamente alterados sem versionamento e cargas que sobrescrevem o histórico. Não há catálogo de dados nem linhagem. Elabore peça técnica com riscos, controles, arquitetura de melhoria e procedimentos de auditoria.
\end{discursiva}

\begin{resumo}[espelho de correção]
Qualidade: validade, unicidade, completude, consistência e temporalidade; histórico: SCD/versionamento, dados imutáveis/raw, CDC; governança: catálogo, owner/steward, linhagem, regras documentadas; auditoria: reconciliação origem-destino, hashes/contagens, amostras, reprocessamento controlado, logs de pipeline; risco: pagamento/decisão com base errada, fraude não detectada e perda de reprodutibilidade.
\end{resumo}

\section{Treino 4 — Peça técnica: IA generativa em processo administrativo}

\begin{discursiva}[peça técnica — até 60 linhas]
Um órgão utiliza solução de IA generativa com RAG para elaborar minutas de parecer. Documentos internos são indexados sem classificação; usuários podem enviar instruções livres; não há validação humana obrigatória; respostas e fontes não são registradas; e o fornecedor usa telemetria do serviço em ambiente externo. Analise riscos e proponha controles técnicos e de governança.
\end{discursiva}

\begin{resumo}[espelho de correção]
Riscos: vazamento, indexação indevida, prompt injection, recuperação não autorizada, alucinação, ausência de accountability e privacidade; controles: classificação/ACL no retrieval, isolamento de tenants, filtro de conteúdo, proteção contra prompt injection, registro de prompt/contexto/fontes/modelo/versão, human-in-the-loop, avaliação de qualidade, DLP, contrato e tratamento de dados, testes adversariais e monitoramento.
\end{resumo}

\section{Treino 5 — Peça técnica: Kubernetes e cadeia de software}

\begin{discursiva}[peça técnica — até 60 linhas]
Uma aplicação pública em Kubernetes utiliza imagens com tag \texttt{latest}, executa contêineres como root, possui secrets em variáveis de ambiente visíveis em manifests e pipeline CI com conta de serviço de privilégios administrativos no cluster. Não há SBOM nem assinatura de imagens. Redija peça técnica com achados, riscos, evidências e recomendações.
\end{discursiva}

\begin{resumo}[espelho de correção]
Cobrar imutabilidade por digest/tag controlada; least privilege, non-root, securityContext; secret manager e criptografia; RBAC restrito; separação CI/CD; assinatura e verificação de artefatos; SBOM e análise de vulnerabilidades; admission policies; logs/auditoria do cluster; evidências de registry, pipeline, RBAC, manifests, runtime e supply chain.
\end{resumo}

\section{Treino 6 — Peça técnica: contratação e sobrepreço}

\begin{discursiva}[peça técnica — até 60 linhas]
Pesquisa de preços para contratação de solução de TI utilizou três cotações de fornecedores, todas aproximadamente 45\% superiores a contratações públicas recentes de objeto comparável. As referências públicas foram descartadas sem justificativa e o quantitativo de licenças foi estimado pelo número total de servidores, embora apenas 38\% utilizem o sistema atual. Analise planejamento, preço, quantitativo e risco ao erário.
\end{discursiva}

\begin{resumo}[espelho de correção]
Questionar comparabilidade, independência e análise crítica das fontes; três cotações não validam preço por si; verificar parâmetros públicos e metodologia; quantitativo deve ter memória de cálculo baseada em demanda; riscos: sobrepreço, superdimensionamento, ociosidade; evidências: ETP/TR, contratos similares, painéis/preços, inventário de usuários ativos, perfis de acesso; recomendar refazer estimativa e dimensionamento.
\end{resumo}

\section{Questões discursivas — 30 linhas}

\begin{discursiva}[Q1 — Zero Trust]
Explique por que Zero Trust não significa ``não confiar em ninguém'' de forma abstrata. Aborde identidade, contexto, privilégio mínimo, verificação contínua e segmentação, e indique três evidências que um auditor buscaria para avaliar sua implementação.
\end{discursiva}
\begin{resumo}[espelho Q1]Definir princípio ``never trust, always verify'' como decisão contextual; autenticação forte, dispositivo, risco, autorização granular, microsegmentação, reavaliação; evidências: políticas IAM/conditional access, MFA, logs de decisões, RBAC/PAM, segmentação e testes.\end{resumo}

\begin{discursiva}[Q2 — CAP e consistência]
Em sistema distribuído sujeito a partições de rede, explique o teorema CAP e por que a escolha arquitetural não pode ser resumida a ``banco NoSQL não tem consistência''. Dê exemplo de trade-off aplicável a sistema público.
\end{discursiva}
\begin{resumo}[espelho Q2]Explicar C, A, P; durante partição existe trade-off entre consistência e disponibilidade; modelos variam por operação/sistema; eventual consistency não é sinônimo de ausência de consistência; exemplo: consulta informativa pode priorizar disponibilidade, enquanto operação financeira pode exigir consistência forte.\end{resumo}

\begin{discursiva}[Q3 — RAG versus fine-tuning]
Compare RAG e fine-tuning em uma solução de IA institucional, considerando atualização de conhecimento, governança da fonte, custo, rastreabilidade e risco de vazamento.
\end{discursiva}
\begin{resumo}[espelho Q3]RAG injeta contexto recuperado e facilita atualização/citação/ACL; fine-tuning altera comportamento/padrões do modelo e não é mecanismo primário para base factual mutável; ambos exigem governança; analisar dados de treinamento, retrieval, acesso, versionamento e avaliação.\end{resumo}

\begin{discursiva}[Q4 — RPO e RTO]
Diferencie RPO e RTO e explique como evidências de backup e testes de restauração permitem verificar se metas declaradas de continuidade são realmente atingíveis.
\end{discursiva}
\begin{resumo}[espelho Q4]RPO = perda temporal tolerável; RTO = tempo para restabelecimento; frequência/replicação sustentam RPO; testes de restauração, capacidade, dependências e runbooks sustentam RTO; política sem teste não demonstra atingimento.\end{resumo}

\begin{discursiva}[Q5 — Evidência de auditoria]
Diferencie suficiência e adequação da evidência de auditoria e explique por que grande quantidade de logs produzidos exclusivamente pela contratada pode continuar sendo evidência fraca.
\end{discursiva}
\begin{resumo}[espelho Q5]Suficiência = quantidade; adequação = qualidade/relevância/confiabilidade; volume não corrige baixa independência, possibilidade de alteração, incompletude ou ausência de cadeia de custódia; buscar fonte independente, integridade, timestamps, controles de acesso e reconciliação.\end{resumo}

\begin{discursiva}[Q6 — LGPD e logs]
Explique como princípios de necessidade, segurança e responsabilização devem orientar retenção de logs que contenham dados pessoais, sem eliminar requisitos de segurança, auditoria e investigação.
\end{discursiva}
\begin{resumo}[espelho Q6]Definir finalidade e base, coletar campos necessários, proteger acesso e integridade, estabelecer retenção justificada, pseudonimizar quando viável, registrar acessos e descarte, preservar logs necessários para obrigação legal/segurança conforme fundamento aplicável.\end{resumo}

\section{Checklist de autocorreção}

Antes de encerrar cada treino, marque:
\begin{itemize}
\item respondi todos os verbos do comando (explique, compare, identifique, proponha)?
\item usei critério técnico/normativo em vez de opinião?
\item conectei condição a risco ou efeito?
\item indiquei evidências verificáveis?
\item minhas recomendações atacam a causa?
\item evitei siglas sem contexto?
\item respeitei o limite de linhas?
\item revisei concordância, regência, pontuação e propriedade vocabular?
\end{itemize}

\section{Plano de treino até a prova}

Realize ao menos uma peça e duas questões por semana. Nas últimas seis semanas, faça blocos completos de 4 horas, com uma peça e duas questões, manuscritos. Registre erros em quatro categorias: \textbf{conteúdo}, \textbf{atendimento ao comando}, \textbf{estrutura} e \textbf{língua portuguesa}. Reescreva as produções em que um subitem essencial foi omitido.
